\chapter{AI Prompt Log}
\ChapterLabel{AIPromptLog}
\section{Chapter Responsibility}
\label{Section:AIPromptLogResponsibility}
Preserve the provenance of material AI assistance and its human review.

\noindent\textbf{Expected result:} An auditable account of how AI contributed to the engineering record.

\section{How to Use This Appendix}
Use this log when AI assistance materially affects requirements, architecture, code, tests, or release evidence. Record enough context for a reviewer to understand the input, output, assumptions, and human corrections.
\section{SP800 Initialization Record}
\label{Section:SP800InitializationLog}
Date: November 14, 2026. Tool: OpenAI Codex; model: GPT-6.
A. Crone instructed the agent to follow course manual version 1.24,
Section 1.6, and confirmed SP800, filename prefix DevFTDI245R, owner
A. Crone, unchanged host-application scope, and Embedded Systems applicability.
Sources: template version 2.33, its bootstrap, AI guide, execution workflow,
document guide, embedded profile, introduction, and course Section 1.6.
One agent performed initialization using
\hyperref[Skill:AIS-01]{AIS-01}, \hyperref[Skill:AIS-02]{AIS-02},
\hyperref[Skill:AIS-13]{AIS-13}, and \hyperref[Skill:AIS-14]{AIS-14}.
The namespace and scope were user decisions; hardware details and numerical
limits were not invented. Template examples remain explicitly illustrative.
The owner approved preparation of this draft; review of the resulting
architecture, baseline acceptance, publication, and release remain pending.
Build evidence is recorded in the project build report; successful compilation
does not establish product verification or validation.

\section{Layout and Navigation Correction}
November 15, 2026. A. Crone authorized correction and rebuild after the first
local review. Codex corrected the context diagram, attribution and profile
links, prompt wrapping, contents-number spacing, and workflow figure size.
The F: project was read-only to this automation session; the corrections and
build were prepared in a separate review copy. The original F: copy remains
unchanged. Owner review and replacement of the working copy remain pending.

\section{Project Metadata Update}
November 15, 2026. A. Crone requested the current project name SP800Project
and organization Stephens University, 2026F SPE 800-ONLB-NGC. The spelling
is retained exactly as supplied. Earlier records retain the previous project
name SP800 as historical provenance. Filename prefix and scope are unchanged.

\section{Example Record}
Date: \textit{[timestamp]}; Tool and model: \textit{[name]}; Input: \textit{[requirements and constraints supplied]}; Task: \textit{[requested output]}; Human review and corrections: \textit{[summary]}; Unresolved assumptions: \textit{[list or none]}.

\section{Requirement Scope or Family Extension Record}
When an AI-assisted change proposes a requirement scope or family that is not
already defined, record the proposed uppercase code, full name, decision
boundary, rationale, representative examples, semantic-overlap review,
identifier-collision check, affected register and category-catalog changes, and
the authorized human's approval or rejection. Do not assign an identifier from
the proposed scope or family before approval.
